\documentclass[11pt, a4paper]{article}

\usepackage{graphicx} % Required for inserting images
% \usepackage[utf8]{inputenc} % Quins caràcters es veuran
\usepackage[document]{ragged2e} % Posar els paràgrafs
\usepackage{amsmath,amsthm,amssymb,latexsym}
\usepackage{enumitem}
% \usepackage{xcolor}
\usepackage[x11names]{xcolor}
\usepackage{tikz}
\usepackage{tcolorbox}
\usepackage{tkz-euclide}
\usepackage{pgfplots}
% \usepackage{cancel}
\usepackage{hyperref}
\usepackage{mathrsfs,mathtools}
\usepackage{microtype}
\usepackage{fancyhdr}
\usepackage{parskip}
\usepackage{tkz-euclide}

\usepackage[catalan]{babel}

\renewcommand\refname{}
\usepackage[margin=1in]{geometry}



\pgfplotsset{compat=1.18}
\usetikzlibrary{patterns}
\usetikzlibrary{intersections}


\setlength{\headheight}{15.21004pt}

\DeclareMathOperator{\Ima}{Im}

%% Macros
\providecommand{\ol}{\overline}
\providecommand{\ul}{\underline}
\providecommand{\wt}{\widetilde}
\providecommand{\wh}{\widehat}
\providecommand{\eps}{\varepsilon}
\providecommand{\half}{\frac{1}{2}}
\providecommand{\inv}{^{-1}}
\newcommand{\dang}{\measuredangle} %% Directed angle
\providecommand{\CC}{\mathbb C}
\providecommand{\FF}{\mathbb F}
\providecommand{\NN}{\mathbb Z_{\geq 1}}
\providecommand{\QQ}{\mathbb Q}
\providecommand{\RR}{\mathbb R}
\providecommand{\PP}{\mathbb P}
\providecommand{\ZZ}{\mathbb Z}
\renewcommand{\AA}{\mathbb A}
\providecommand{\dd}{\mathrm d}
\newcommand{\mfrk}{\mathfrak}
\newcommand{\KK}{\mathbb{K}}
\newcommand{\TT}{\mathcal{T}}
\providecommand{\ts}{\textsuperscript}
\providecommand{\dg}{^\circ}
\providecommand{\ii}{\item}
\DeclareMathOperator*{\lcm}{lcm}
\DeclareMathOperator*{\argmin}{arg min}
\DeclareMathOperator*{\argmax}{arg max}
\DeclareMathOperator*{\rad}{rad}
\DeclareMathOperator*{\spec}{Spec}
\DeclareMathOperator*{\mult}{mult}
\providecommand{\bigO}{\mathcal O}




% theorem environments
% Les versions estrellades no estan numerades 
\theoremstyle{definition}
    \newtheorem{theorem}{Teorema}
    \newtheorem{lemma}[theorem]{Lema}
    \newtheorem{claim}[theorem]{Claim}
    \newtheorem{corollary}[theorem]{Corol·lari}
    \newtheorem{definition}[theorem]{Definició}
    \newtheorem{proposition}[theorem]{Proposició}
    \newtheorem{example}[theorem]{Exemple}
    
    \newtheorem*{theorem*}{Teorema}
    \newtheorem*{lemma*}{Lema}
    \newtheorem*{claim*}{Claim}
    \newtheorem*{corollary*}{corol·lari}
    \newtheorem*{definition*}{Definició}
    \newtheorem*{proposition*}{Proposició}
    \newtheorem*{example*}{Exemple}
    \newtheorem*{proof*}{Demostració}
    
\theoremstyle{remark}
    \newtheorem{remark}[theorem]{Remarca}
    \newtheorem*{remark*}{Remarca}

\usepackage{hyperref}
\renewcommand{\thesubsection}{}
\renewcommand{\thesection}{}

\begin{document}
\pagestyle{fancy}
% 
\lhead{Bernat Esteve}
\rhead{Parcial 2024}
\begin{tcolorbox}[title=\section{Exercici 1}(2 punts)]
Calculeu raonadament la multiplicitat d'intersecció en el punt $(0, 0)$ dels parells de corbes següents:
\begin{enumerate}
    \item $x^3 + y^4$ i $x^4 + y^3$.
    \item $y+x^2+y^2$ i $y^3+x^4+y^4$.
    \item $x^2y^2$ i $x^2+y^2+x^3+y^3$
\end{enumerate}
\end{tcolorbox}
\subsection{Exercici 1.1}
Notem que el con tangent de la primera corba és $x=0$, i el de la segona és $y=0$, per tant, tenim que la multiplicitat és $3\times 3=9$.
\subsection{Exercici 1.2}
Notem que el con tangent de la primera corba i el de la segona són iguals ($y=0$). Per tant, hem de fer alguna maniobra:
\[
    y^3+x^4+y^4-y^2(y+x^2+y^2)=x^4-x^2y^2
\]
Que no té $y=0$ en el con tangent, per tant, per l'axioma de l'ideal en la definició de multiplicitat d'intersecció, ens queda que la multiplicitat és $1\times 4=4$.
\subsection{Exercici 1.3}
Notem que el con tangent de la primera corba és $x=0$ i $y=0$; en canvi a la segona és $x+iy=0$ i $x-iy=0$; i per tant, tindríem que la multiplicitat és $4\times 2=8$.
\newpage
\begin{tcolorbox}[title=\section{Exercici 2} (2 punts)]
    Considerem les còniques projectives $C_1 = V_+(x^2-yz)$, $C_2=V_+(x^2-xy-yz)$, $C_3=V_+(yz-x^2+y^2)$.
    \begin{enumerate}
        \item Proveu que són no singulars.
        \item Trobeu la intersecció de $C_1$ i $C_2$ i calculeu la multiplicitat d'intersecció en cada punt de la intersecció.
        \item Feu el mateix que en l'apartat b) per a $C_1$ i $C_3$.
    \end{enumerate}
\end{tcolorbox}
\subsection{Exercici 2.1}
Mirant $C_1$
\[
    \begin{cases}
        \partial_x C_1=2x=0\\
        \partial_y C_1=-z=0\\
        \partial_z C_1=-y=0\\
    \end{cases}
\]
Que no té solució. Fent el mateix amb les altres corbes:
\[
    \begin{cases}
        \partial_x C_2=2x-y=0\\
        \partial_y C_2=-x-z=0\\
        \partial_z C_2=-y=0\\
    \end{cases}
    \qquad
    \begin{cases}
        \partial_x C_3=-2x=0\\
        \partial_y C_3=2y+z=0\\
        \partial_z C_3=y=0\\
    \end{cases}
\]
Que ens diu que cap d'elles té cap punt singular.
\subsection{Exercici 2.2}
Considerem el sistema d'equacions:
\[
    \begin{cases*}
        x^2-yz=0\\
        x^2-yz-xy=0
    \end{cases*}\implies
    \begin{cases*}
        x^2-yz=0\\
        xy=0
    \end{cases*}
\]
Per tant, hem de considerar 2 casos:\par
\textbf{Cas 1: $y=0$}\\
Si $y=0$, mirant a les equacions, ens dona $x=0$. Per tant, un punt és $(0:0:1)$ (que efectivament és punt).\par
\textbf{Cas 2: $z=0$}\\
Mirant les equacions, això també ens dona $x=0$, per tant, tenim $(0:1:0)$.\par

Mirem ara les multiplicitats d'intersecció. Considerem la recta $r_\infty=\{y=0\}$. Aleshores
\[
I_{(0,0)}(x^2-z,x^2-x-z)=
I_{(0,0)}(x^2-z,-x)=
I_{(0,0)}(-z,-x)=1
\]
Per tant, L'altre punt té intersecció 3.
\subsection{Exercici 2.3}
Fem el mateix d'abans. Considerem el sistema d'equacions:
\[
    \begin{cases*}
        x^2-yz=0\\
        x^2-y^2-zy=0
    \end{cases*}\implies
    \begin{cases*}
        x^2-yz=0\\
        y^2=0
    \end{cases*}
\]
Per tant, a diferència d'abans, només hem de considerar 1 sol cas:\par
Si $y=0$, mirant a les equacions, ens dona $x=0$. Per tant, un punt és $(0:0:1)$. I com que és únic, ha de tenir multiplicitat 4.
\newpage
\begin{tcolorbox}[title=\section{Exercici 3}]
Trobeu l'equació de la corba que satisfan els punts de la forma $(t^3 + t, t^2 - 1)$, on $t \in \CC$. Doneu les singularitats d'aquesta corba.
\end{tcolorbox}
\subsection{Exercici 3}
Notem que 
\[
    x=t(t^2+1)=t(y+2)
\]
Per tant:
\[
    x^2-(y+2)^2=(t^2-1)(y+2)^2=(y+2)^2y
\]
És a dir, que $\gamma$ satisfà:
\[
    x^2=(y+2)^2(y+1)
\]
Per veure l'altra inclusió, sigui $(x,y)$ que satisfà aquesta equació. Aleshores sigui $t$ tal que $x=t(y+2)$:
\[
    x^2=t^2(y+2)^2=(y+2)^2(y+1)
\]
Per tant
\[
    (t^2-1-y)(y+2)^2=0
\]
Que ens diu que $y+2=0$, i mirant l'equació ens dona $x=0$; o que $y=t^2-1$.\par
I si tenim que $y=t^2-1$, aleshores, posant això a l'equació inicial, ja estem.\par
Per mirar els punts singulars, hem de mirar les derivades parcials de la corba:

\[
    \begin{cases*}
    f=x^2-(y+2)^2(y+1)=0\\
    \frac{\partial f}{\partial x}=2x=0\\
    \frac{\partial f}{\partial y}=(y+2)(3y+4)=0\\
    \end{cases*}
\]
Per tant, tenim que $x=0$, i que $y=-2$ o $y=-\frac{4}{3}$. Posant aquests punts a la corba, només el $(0,-2)$ és punt singular.\par
Per mirar la multiplicitat, centrem el polinomi en $(0,-2)$:
\[
    x^2-(y+2)^2(y+1)\to \ol x^2 - \ol y^2-\ol y^3
\]
Per tant, té multiplicitat 2.\par \vspace{20mm}
La manera oficial de fer aquest problema consisteix en utilitzar la resultant. Recordem que hem de pensar en la resultant com una funció de $\CC[x_1,\dots,x_n]^2\to \CC[x_2,\dots, x_n]$ que val zero en $(x_2,\dots, x_n)$ si i només si $F$ i $G$ s'intersequen en $(x_1,\dots,x_n)$ per algun $x_1$. Per tant, si considerem:
\[
    \begin{cases}
        F=x(t)-x\\
        G=y(t)-y
    \end{cases}
\]
Que fixada una $(x,y)$ donarà zero si i només $\exists t$ tal que $x(t)=x$ i $y(t)$. Per tant, si fem $R(x(t)-x,y(t)-y;t)$:
\[
    \det\begin{pmatrix}
        1 & 0 &    1 &   -x & 0\\
        0 & 1 &    0 &    1 & -x\\
        1 & 0 & -1-y &    0 & 0\\
        0 & 1 &    0 & -1-y & 0\\
        0 & 0 &    1 &    0 & -1-y\\
    \end{pmatrix}=x^2-4-8y-5y^2-y^3=0
    % -2-3y-y^2+x^2-1-3y-3y^2-y^3-1-2y-y^2
\]
Que dona el mateix que ens havia donat abans.
\newpage
\begin{tcolorbox}[title=\section{Exercici 4}]
    Considerem la cúbica del pla projectiu d'equació: $-xy^2 + xyz + (z - x)^3=0$.\\
Demostreu que és no singular i que té una inflexió en el punt $[1 : 1 : 1]$.
\end{tcolorbox}
\subsection{Exercici 4}
Calculem les derivades:
\[
    \begin{cases}
        \partial_xF=-y^2+yz-3(z-x)^2=0\\
        \partial_yF=x(-2y+z)=0\\
        \partial_zF=yx+3(z-x)^2=0\\
    \end{cases}
\]
Per tant, tenim 2 opcions: $x=0$ o $z=2y$.\\
\textbf{Cas 1: $x=0$}\\
Posant $x=0$ a les equacions, obtenim:
\[
    \begin{cases}
        -y^2+yz-3z^2=0\\
        3z^2=0\\
    \end{cases}\implies
    \begin{cases}
        z=0\\
        y=0\\
    \end{cases}
\]
Que és una contradicció.\\
\textbf{Cas 2: $y=2z$}\\
\[
    \begin{cases}
        -2z^2-3(z-x)^2=0\\
        2zx+3(z-x)^2=0\\
    \end{cases}
\]
Sumant les equacions tenim:
\[
    \begin{cases}
        -2z^2-3(z-x)^2=0\\
        z(x-z)=0\\
    \end{cases}
\]
Per tant, tenim $z=0$, que ens dona $x=0$; per tant, no pot ser. I si fem $x=z$ ens torna a donar que $z=0$.\\
Per tant, aquesta cúbica no té punts singulars.
\end{document}
% ---
% \newpage
% \begin{tcolorbox}[title=\section{Exercici XXX}]
% \end{tcolorbox}
% \subsection{Exercici XXX}
